\section*{\centering Resumo}
\begin{resumotexto}
Este documento demonstra uma organização LaTeX para artigo científico em português brasileiro, tomando como exemplo a relação entre um documento discente de design de jogo e o código do Night Heist. O objetivo é apresentar uma estrutura editável que separe problema, objetivo, método, evidência, interpretação e limites do estudo. A demonstração utiliza informações já registradas no projeto e não produz novos dados, medições ou resultados educacionais. O método exemplificado é a análise documental com inspeção estática, baseada em localizadores da especificação, critérios observáveis, símbolos do código e estado de revisão. Um quadro mostra três situações presentes na matriz piloto: correspondência estrutural, divergência literal e impossibilidade de comparação por ausência de parâmetro. A proposta editorial aplica resumo em parágrafo único, palavras-chave padronizadas, seções numeradas, quadro com fonte e referências autor-data. O exemplo não demonstra conformidade integral com todas as normas vigentes e não constitui template oficial de periódico. As classificações técnicas utilizadas permanecem preliminares e dependem de revisão humana. A contribuição demonstrativa é uma base reutilizável que permite substituir conteúdo e metadados sem alterar os comandos de apresentação \citep{corpus2026,abnt2021}.
\end{resumotexto}
\noindent\textbf{Palavras-chave:} game design; rastreabilidade; análise documental; LaTeX; artigo científico.

\noindent\begin{minipage}{\linewidth}
\section*{\centering Abstract}
\begin{resumotexto}
This document demonstrates an editable LaTeX structure for a scientific article in Brazilian Portuguese, using the relationship between a student game design document and the Night Heist code as an example. Its objective is to separate the research problem, objective, method, evidence, interpretation, and limitations. The demonstration reuses information already recorded in the project and produces no new measurements or educational findings. The illustrated method is documentary analysis with static inspection, based on specification locators, observable criteria, code symbols, and review status. A qualitative comparison presents three situations from the pilot matrix: structural correspondence, literal divergence, and inability to compare an unspecified parameter. The editorial proposal includes a single-paragraph abstract, standardized keywords, numbered sections, source-labelled qualitative comparisons, and author-date references. It does not establish full compliance with all current standards and is not an official journal template. The technical classifications remain preliminary and require human review. The demonstration contributes a reusable structure in which content and metadata can be replaced independently from presentation commands \citep{corpus2026,abnt2021}.
\end{resumotexto}
\noindent\textbf{Keywords:} game design; traceability; documentary analysis; LaTeX; scientific article.
\end{minipage}
